\chapter{The Idea Monoid}\label{ch:monoid}

\emph{In which the bare combinatorial skeleton of Idea Theory is laid down.
We isolate three axioms --- associativity and a two-sided unit --- which we
defend on independent philosophical grounds; we show that they suffice to
derive a full apparatus of powers, chains, sub-ideas, idempotents,
cancellation, and invertibility; and we prove the three headline theorems
(\ref{thm:1-1}--\ref{thm:1-3}) and four corollaries
(\ref{cor:1-1}--\ref{cor:1-4}) that organize the rest of the volume.  Every
result of this chapter is machine-verified in the Lean~4 module
\texttt{IdeaTheory.Theorems1}; we cite the verified identifier in
parentheses.  The chapter closes by showing that the bare monoid already
suffices to host four classical positions on the nature of ideas: those of
\cite{Locke1690}, \cite{Hume1739}, \cite{Frege1892}, and the later
\cite{Wittgenstein1953}.}

\section{The axioms of an idea algebra}\label{sec:axioms}

The whole edifice of Idea Theory rests on three axioms.  We isolate them
here, before introducing any further analytic structure, in order to make
clear which results require nothing more than the bare monoid skeleton.

\begin{definition}[Idea algebra; bare layer]\label{def:idea-algebra}
A \emph{bare idea algebra} is a triple $\IdMon$ in which $\Ideas$ is a set
(the set of \emph{ideas}), $\compose : \Ideas \times \Ideas \to \Ideas$ is
a binary operation (\emph{composition}), and $\unit \in \Ideas$ is a
distinguished element (the \emph{void idea}, or unit), subject to the
following three axioms.
\begin{description}
\item[(A1)\ Associativity.]
$(a \compose b) \compose c \;=\; a \compose (b \compose c)$ for all
$a,b,c \in \Ideas$.
\item[(A2)\ Left unit.]
$\unit \compose a \;=\; a$ for all $a \in \Ideas$.
\item[(A3)\ Right unit.]
$a \compose \unit \;=\; a$ for all $a \in \Ideas$.
\end{description}
\end{definition}

The mathematical content of Definition~\ref{def:idea-algebra} is exactly that
of a monoid in the sense of \cite{MacLane1998} or \cite{Bourbaki1989}; what
is new is the interpretation we attach to it.  A few words on each axiom
are in order, since the philosophical literature on ideas is the only
arena in which we must be prepared to defend any of them.

\paragraph{Why associativity (A1).}
Associativity is the algebraic shadow of \emph{bracket-independence of
content}.  When two ideas $a$ and $b$ combine, the result $a \compose b$ is
itself an idea; if a third idea $c$ is then brought into contact, we obtain
an idea $(a \compose b) \compose c$.  Yet the same three ideas could
equally have been combined in the order $a$, then $b \compose c$,
producing $a \compose (b \compose c)$.  Axiom (A1) declares that these two
ideas are identical, not merely equivalent.

The empirical pull behind (A1) is strong.  Hume's principle of
association (\cite{Hume1739}, Bk.~I, Pt.~I, \S IV) is most naturally read
as a non-bracketed enumeration: ``A succeeded by B succeeded by C''.
Frege's compositionality (\cite{Frege1892}) requires that the sense of a
complex expression depend only on the senses of its parts and their
arrangement; bracketing is not part of the arrangement once the linear
order has been fixed.  Whitehead's process metaphysics, Hegel's repeated
\emph{Aufhebungen}, and Sperber's chains of cultural attractors all
presuppose that bracketing is invisible at the level of content.  We
therefore include (A1) without further justification, treating it as a
\emph{minimal} commitment that any honest theory of combinable ideas must
make.

We do \emph{not} assume commutativity.  ``Newtonian mechanics composed
with quantum mechanics'' is a different idea from ``quantum mechanics
composed with Newtonian mechanics'', and the burden of evidence is on the
philosopher who would equate them.  This non-commutativity will be
sharpened in Chapter~2, where the antisymmetric part of the resonance
pairing measures exactly the failure of $\compose$ to be commutative.

\paragraph{Why a two-sided unit (A2--A3).}
The void idea $\unit$ models the \emph{empty thought} --- the contentless
idea that, when prepended or appended to any other, leaves it unchanged.
Several distinct philosophical considerations converge on the existence of
$\unit$.

First, every model we shall consider possesses a natural ``do nothing''
element: in Locke's combinatorial picture (\cite{Locke1690}, Bk.~II,
Ch.~XII), it is the trivial complex idea built from no simple ideas at
all; in modern semantic-vector models it is the zero context vector; in
narrative theory it is the empty episode.  Second, technical considerations
favor the unit: without it the powers $\compose^n$ have no natural value
at $n=0$, the empty chain has no natural composite, and the cohomology
class of the emergence cocycle in Chapter~5 is no longer well defined.
Third, the unit is harmless: if a candidate carrier $\Ideas'$ has no unit
we may freely adjoin one, defining $\Ideas \defeq \Ideas' \sqcup \{\unit\}$
with $\unit \compose a \defeq a$ and $a \compose \unit \defeq a$.

We split the unit axiom into a left and a right form (A2, A3).  Either
half could in principle hold without the other, and we shall see in
\S\ref{sec:uniqueness} that requiring both is strictly stronger than
requiring either one in isolation.

\paragraph{What we are \emph{not} assuming.}
It is worth dwelling on the negative content of
Definition~\ref{def:idea-algebra}.  We do not assume that $\compose$ is
\emph{commutative}: as already remarked, the order in which two ideas
combine is in general semantically loaded.  We do not assume that
$\compose$ has \emph{absorbing} elements: there is no a~priori candidate
for a ``zero idea'' that swallows every composition (although particular
models, such as the saturation algebra of Example~\ref{ex:two-atom}
below, may exhibit one).  We do not assume \emph{idempotency} of any idea
other than $\unit$ (and indeed the idempotency of $\unit$ is itself a
\emph{theorem}, Lemma~\ref{lem:unit-self}, not an axiom).  We do not
assume that ideas can be \emph{cancelled} (we do not require $a \compose
b = a \compose c \Rightarrow b = c$); whether a given idea is
cancellative is a definable property treated in
\S\ref{sec:cancellation}, not a structural commitment.  Finally, we do
not assume any \emph{order}, \emph{topology}, or \emph{measure} on
$\Ideas$; the bare layer is purely combinatorial.

This austerity is deliberate.  Every additional axiom one builds in is an
exegetical decision that prejudges some debate in the literature.  The
philosophical purchase of Idea Theory comes from \emph{deriving} as much
as possible from the minimal axiom set, and reserving the right to add
further structure only when the empirical question being modelled
demands it.  Chapters 2--9 of Volume~I will progressively layer more
structure onto $\IdMon$; the discipline of starting from three axioms
guarantees that, when those layers do their work, we know exactly which
of our conclusions depend on which structural assumptions.

\begin{remark}[Why \emph{three} axioms, not more]
We deliberately resist the temptation to package commutativity, idempotency,
or cancellation into the bare definition.  Each of these is a substantive
claim about a specific class of ideas (cf.\ \S\ref{sec:idempotents}--%
\ref{sec:cancellation}), and conflating them with the structural skeleton
would force the formalism to take sides on disputed questions before the
relevant phenomena have been described.  Three axioms are exactly what is
needed to organize the literature, no more.
\end{remark}

\begin{remark}[On the choice of name ``idea algebra'']
We refer to the structure $\IdMon$ as an idea \emph{algebra} rather than
an idea \emph{monoid} for two reasons.  The first is anticipatory: the
later chapters extend $\IdMon$ with additional algebraic data --- a
bilinear pairing, a graded decomposition, a 2-cocycle, a grading --- and
the resulting object is no longer a monoid in any standard sense.  The
second is exegetical: the philosophical literature uses the word
``algebra'' loosely, and our use is intended to evoke the entire stack
of structure that culminates in the full Definition of an idea algebra
in the introduction (\S 1 of the brief).  In this chapter the word
``algebra'' should be read as ``algebra-with-only-the-bare-monoid-layer''.
The risk of confusion is small in practice because every chapter declares
which layers are in scope at the start.
\end{remark}

\begin{notation}[Operator conventions]
Throughout this chapter the symbols $a, b, c, x, y, z, p, q$ denote
arbitrary elements of $\Ideas$; $m, n, k$ denote natural numbers;
$\unit$ is always the unit of $\IdMon$.  We write $a \compose b \compose c$
for $(a \compose b) \compose c$ (which equals $a \compose (b \compose c)$
by (A1)), and we write $a^n$ for the $n$-fold composition $\npow(n,a)$
defined in \S\ref{sec:powers}.  Functions and relations introduced in this
chapter are summarized in the table at the end of \S\ref{sec:cancellation}.
\end{notation}

\section{Uniqueness of the unit and basic congruence lemmas}\label{sec:uniqueness}

The first non-trivial fact about $\IdMon$ is that the void idea is
\emph{unique} as a two-sided unit.  This is the standard monoid argument,
included here both for completeness and because it is the prototype for
many uniqueness arguments in the chapters that follow.

\begin{lemma}[Uniqueness of two-sided units]\label{lem:unit-unique-twosided}
Let $e \in \Ideas$ satisfy $e \compose a = a$ and $a \compose e = a$ for
every $a \in \Ideas$.  Then $e = \unit$.
\end{lemma}

\begin{proof}
By the left-unit hypothesis applied at $a = \unit$ we have
$e \compose \unit = \unit$.  By (A3) applied to $e$ we have
$e \compose \unit = e$.  Equating the two right-hand sides yields
$e = \unit$.  (Verified as
\texttt{IdeaTheory.ident\_unique\_two\_sided}.)
\end{proof}

The same argument can be sharpened: a one-sided unit is already enough,
provided one accepts the corresponding axiom on the other side.

\begin{lemma}[Left units are units]\label{lem:left-unit}
If $e \in \Ideas$ satisfies $e \compose a = a$ for every $a$, then
$e = \unit$.  Symmetrically, if $a \compose e = a$ for every $a$, then
$e = \unit$.
\end{lemma}

\begin{proof}
For the first claim, take $a = \unit$ in the hypothesis to get
$e \compose \unit = \unit$.  By (A3), $e \compose \unit = e$, so $e = \unit$.
The second claim is dual: take $a = \unit$ to get $\unit \compose e = \unit$,
hence by (A2) $e = \unit$.  (Verified as
\texttt{IdeaTheory.left\_ident\_eq\_ident} and
\texttt{IdeaTheory.right\_ident\_eq\_ident}.)
\end{proof}

\begin{remark}
Lemma~\ref{lem:left-unit} shows that the apparent redundancy of axioms
(A2) and (A3) is only apparent: each of them, \emph{together with the
other}, suffices to pin down $\unit$ uniquely.  Without one of the two,
there can be many one-sided units (e.g.\ in a left-zero band every element
is a left unit).
\end{remark}

We next collect three congruence lemmas which formalize the principle
that composition respects equality of arguments.  These are completely
trivial mathematically but they will be invoked dozens of times throughout
the book, and naming them once spares us repetition later.

\begin{lemma}[Left congruence]\label{lem:left-cong}
If $b = c$, then $a \compose b = a \compose c$ for every $a$.
\end{lemma}

\begin{lemma}[Right congruence]\label{lem:right-cong}
If $b = c$, then $b \compose a = c \compose a$ for every $a$.
\end{lemma}

\begin{lemma}[Bilinear congruence]\label{lem:bilin-cong}
If $a = c$ and $b = d$, then $a \compose b = c \compose d$.
\end{lemma}

\begin{proof}
All three lemmas are immediate from the substitution principle: if two
arguments are equal, replacing one by the other inside any expression
preserves the value.  (Verified as \texttt{IdeaTheory.op\_left\_congr},
\texttt{op\_right\_congr}, \texttt{op\_congr}.)
\end{proof}

The role of these elementary lemmas is foundational: every later
manipulation of equations involving $\compose$ implicitly uses one of
them.  In a setting where $\compose$ were a partial operation, or where
equality of ideas were not assumed to be a congruence, one of the three
congruence lemmas would fail and the entire algebraic apparatus would
collapse.  We record them now so that the reader can verify, when in
doubt, that no further hidden axiom is being smuggled in.

\begin{lemma}[Stability of $\unit$ under self-composition]\label{lem:unit-self}
$\unit \compose \unit = \unit$.
\end{lemma}

\begin{proof}
Take $a = \unit$ in (A2).  (Verified as
\texttt{IdeaTheory.op\_ident\_ident}.)
\end{proof}

Lemma~\ref{lem:unit-self} foreshadows the more general fact, proved in
\S\ref{sec:idempotents}, that $\unit$ is the canonical example of an
\emph{idempotent} idea --- one which is unaffected by being composed with
itself.

\section{Powers of an idea}\label{sec:powers}

The first construction made possible by the bare monoid axioms is
\emph{iteration}.  Informally, we wish to make precise what it means to
``invoke the same idea $n$ times in succession''.  This is the formal
correlate of repeated metaphor (\cite{Hume1739}, on the strengthening of
association by repetition), liturgical recitation, repeated theorem
application, or repeated activation of a neural ensemble.

\begin{definition}[Iterated composition]\label{def:npow}
The \emph{$n$-th power} of an idea $a \in \Ideas$, written $a^n$ or
$\npow(n,a)$, is defined by recursion on $n$:
\begin{equation}\label{eq:npow-def}
a^0 \defeq \unit,
\qquad
a^{n+1} \defeq a \compose a^n.
\end{equation}
\end{definition}

\begin{remark}
The convention $a^0 = \unit$ is forced upon us once we accept the unit
axiom and the additive law (Lemma~\ref{lem:npow-add} below).  Without
$\unit$ the value of $a^0$ would be undefined and the additive law would
acquire ugly side conditions.
\end{remark}

\begin{lemma}[Elementary identities for $\npow$]\label{lem:npow-elem}
For every $a \in \Ideas$ and $n \in \NN$,
\begin{align}
a^0 &= \unit,
\label{eq:npow-zero} \\
a^{n+1} &= a \compose a^n,
\label{eq:npow-succ-left} \\
a^1 &= a,
\label{eq:npow-one} \\
\unit^n &= \unit.
\label{eq:npow-unit}
\end{align}
\end{lemma}

\begin{proof}
\eqref{eq:npow-zero} and \eqref{eq:npow-succ-left} are immediate from
Definition~\ref{def:npow}.  For \eqref{eq:npow-one}, expand:
$a^1 = a \compose a^0 = a \compose \unit = a$ by (A3).  For
\eqref{eq:npow-unit}, induct on $n$: the base case $n = 0$ is
\eqref{eq:npow-zero}; for the inductive step, $\unit^{n+1} = \unit \compose
\unit^n = \unit$ by (A2) and the inductive hypothesis.
(Verified as \texttt{IdeaTheory.npow\_zero}, \texttt{npow\_succ},
\texttt{npow\_one}, \texttt{npow\_ident}.)
\end{proof}

The recursive definition \eqref{eq:npow-def} unwraps the new factor on
the \emph{left}.  The next lemma is the dual: the new factor can equally
be unwrapped on the right.

\begin{lemma}[Right-unwinding of powers]\label{lem:npow-succ-right}
$a^{n+1} = a^n \compose a$ for every $a \in \Ideas$ and $n \in \NN$.
\end{lemma}

\begin{proof}
We argue by induction on $n$.  For $n = 0$,
\[
a^1 \;=\; a \compose a^0 \;=\; a \compose \unit \;=\; a
\;=\; \unit \compose a \;=\; a^0 \compose a.
\]
For the inductive step, assume the claim for $n = k$.  Then
\[
a^{k+2}
\;=\; a \compose a^{k+1}
\;=\; a \compose (a^k \compose a)
\;=\; (a \compose a^k) \compose a
\;=\; a^{k+1} \compose a,
\]
where the second equality is the inductive hypothesis and the third is
(A1).  (Verified as \texttt{IdeaTheory.npow\_succ'}.)
\end{proof}

\begin{lemma}[Self-commutativity of powers]\label{lem:npow-self-commute}
For every $a$ and $k$, $a \compose a^k = a^k \compose a$.
\end{lemma}

\begin{proof}
Both expressions equal $a^{k+1}$: the left-hand side by definition, the
right-hand side by Lemma~\ref{lem:npow-succ-right}.  (Verified as
\texttt{IdeaTheory.npow\_comm\_self}.)
\end{proof}

The two main arithmetic laws of iteration --- additivity in the exponent
and multiplicativity in the exponent --- now follow.

\begin{lemma}[Additivity in the exponent]\label{lem:npow-add}
For every $a \in \Ideas$ and every $m, n \in \NN$,
\begin{equation}\label{eq:npow-add}
a^{m+n} \;=\; a^m \compose a^n.
\end{equation}
\end{lemma}

\begin{proof}
Induct on $m$.  For $m = 0$ both sides equal $a^n$: the left-hand side
by \eqref{eq:npow-zero}'s consequence $0+n = n$, the right-hand side by
\eqref{eq:npow-zero} and (A2).  Assume the claim for $m = k$.  Then
\begin{align*}
a^{(k+1)+n}
&= a^{(k+n)+1}
= a \compose a^{k+n} \\
&= a \compose (a^k \compose a^n)
= (a \compose a^k) \compose a^n
= a^{k+1} \compose a^n,
\end{align*}
using the inductive hypothesis at the third step and (A1) at the fourth.
(Verified as \texttt{IdeaTheory.npow\_add}.)
\end{proof}

\begin{lemma}[Multiplicativity in the exponent]\label{lem:npow-mul}
For every $a \in \Ideas$ and every $m, n \in \NN$,
\begin{equation}\label{eq:npow-mul}
a^{mn} \;=\; (a^n)^m.
\end{equation}
\end{lemma}

\begin{proof}
Induct on $m$.  For $m = 0$: $(a^n)^0 = \unit = a^0 = a^{0\cdot n}$.  Assume
the claim at $m = k$.  Then $(k+1)n = n + kn$, so
\begin{align*}
a^{(k+1)n}
&= a^{n + kn}
= a^n \compose a^{kn} \\
&= a^n \compose (a^n)^k
= (a^n)^{k+1},
\end{align*}
using Lemma~\ref{lem:npow-add} at the second step, the inductive
hypothesis at the third, and Definition~\ref{def:npow} at the fourth.
(Verified as \texttt{IdeaTheory.npow\_mul}.)
\end{proof}

Together, Lemmas~\ref{lem:npow-add} and~\ref{lem:npow-mul} say that, for
each fixed idea $a$, the function $n \mapsto a^n$ is a homomorphism from
the additive monoid $(\NN, +, 0)$ into $\IdMon$, and the function
$(m,n) \mapsto a^{mn}$ factors through this homomorphism in the
expected way.  These two facts will be re-stated and packaged as
Theorem~\ref{thm:1-3} below.

\begin{lemma}[Small powers]\label{lem:npow-small}
$a^2 = a \compose a$ and $a^3 = a \compose (a \compose a)$.
\end{lemma}

\begin{proof}
$a^2 = a \compose a^1 = a \compose a$ by Lemma~\ref{lem:npow-elem};
$a^3 = a \compose a^2 = a \compose (a \compose a)$.  (Verified as
\texttt{IdeaTheory.npow\_two}, \texttt{npow\_three}.)
\end{proof}

\section{Chains}\label{sec:chains}

Iteration treats the same idea applied $n$ times.  More generally we wish
to compose a finite sequence of (possibly distinct) ideas.  The resulting
construction --- the \emph{chain} of a list of ideas --- is the bridge
between the algebraic apparatus and the empirical phenomena of discourse,
narrative, and cultural transmission.

\begin{definition}[Chain composite]\label{def:chain}
For a finite list $L = [a_1, a_2, \dots, a_k]$ of ideas, the
\emph{chain composite} $\chain{L}$ is defined recursively by
\begin{equation}\label{eq:chain-def}
\chain{[\,]} \defeq \unit,
\qquad
\chain{a :: L} \defeq a \compose \chain{L}.
\end{equation}
The \emph{length} of $L$, written $|L|$, is its number of entries.  The
\emph{void chain} of length $n$ is the list $\unit^{::n}$ consisting of
$n$ copies of $\unit$.
\end{definition}

The chain construction reduces to the more familiar operations as
expected: a chain of length zero is the void idea, a chain of length one
is the single idea itself, and a chain consisting entirely of repeated
copies of $a$ is the corresponding power.

\begin{lemma}[Singleton and constant chains]\label{lem:chain-basic}
For every $a \in \Ideas$ and $n \in \NN$:
\begin{align}
\chain{[a]} &= a,
\label{eq:chain-singleton} \\
\chain{[\underbrace{a,\dots,a}_{n}]} &= a^n,
\label{eq:chain-replicate} \\
\chain{[\underbrace{\unit,\dots,\unit}_{n}]} &= \unit.
\label{eq:chain-void}
\end{align}
\end{lemma}

\begin{proof}
\eqref{eq:chain-singleton}: by Definition~\ref{def:chain},
$\chain{[a]} = a \compose \chain{[\,]} = a \compose \unit = a$ by (A3).
\eqref{eq:chain-replicate}: induct on $n$.  Base case $n=0$: both sides
equal $\unit$.  Inductive step: the chain on the left expands to
$a \compose \chain{[a,\dots,a]_{n}} = a \compose a^n = a^{n+1}$ by
the inductive hypothesis and Definition~\ref{def:npow}.
\eqref{eq:chain-void}: a special case of \eqref{eq:chain-replicate} with
$a = \unit$, combined with Lemma~\ref{lem:npow-elem}\eqref{eq:npow-unit}.
(Verified as \texttt{IdeaTheory.chain\_singleton},
\texttt{chain\_replicate}, \texttt{chain\_replicate\_ident}.)
\end{proof}

The most useful structural fact about chains is the following splitting
identity, which exhibits $\chain{\,\cdot\,}$ as a homomorphism from the
free monoid of lists into $\IdMon$.

\begin{lemma}[Concatenation of chains]\label{lem:chain-append}
For any two finite lists $L_1, L_2$ of ideas,
\begin{equation}\label{eq:chain-append}
\chain{L_1 \mathbin{+\!+} L_2} \;=\; \chain{L_1} \compose \chain{L_2}.
\end{equation}
\end{lemma}

\begin{proof}
By induction on $L_1$.  If $L_1 = [\,]$, both sides equal $\chain{L_2}$:
the left by definition of list concatenation, the right by (A2) and
$\chain{[\,]} = \unit$.  Suppose the identity holds for $L_1 = L$, and let
$L_1 = a :: L$.  Then
\begin{align*}
\chain{(a :: L) \mathbin{+\!+} L_2}
&= \chain{a :: (L \mathbin{+\!+} L_2)} \\
&= a \compose \chain{L \mathbin{+\!+} L_2} \\
&= a \compose (\chain{L} \compose \chain{L_2}) \\
&= (a \compose \chain{L}) \compose \chain{L_2}
 = \chain{a :: L} \compose \chain{L_2},
\end{align*}
using the inductive hypothesis at the third step and (A1) at the fourth.
(Verified as \texttt{IdeaTheory.chain\_append}.)
\end{proof}

\begin{lemma}[Splitting at an arbitrary index]\label{lem:chain-split}
For every list $L$ and every $n \in \NN$,
\begin{equation}\label{eq:chain-split}
\chain{L} \;=\; \chain{\mathrm{take}_n L} \compose \chain{\mathrm{drop}_n L},
\end{equation}
where $\mathrm{take}_n L$ is the prefix of length $\min(n, |L|)$ and
$\mathrm{drop}_n L$ is the corresponding suffix.
\end{lemma}

\begin{proof}
We have $L = (\mathrm{take}_n L) \mathbin{+\!+} (\mathrm{drop}_n L)$ by the
defining property of \texttt{take} and \texttt{drop} for lists.
Substituting into Lemma~\ref{lem:chain-append} yields the claim.
(Verified as \texttt{IdeaTheory.chain\_split}; this is the engine of
Theorem~\ref{thm:1-1} below.)
\end{proof}

\begin{lemma}[Identity-padding]\label{lem:chain-pad}
Let $L$ be any list and let $L'$ be obtained from $L$ by inserting one
copy of $\unit$ at any position.  Then $\chain{L'} = \chain{L}$.
\end{lemma}

\begin{proof}
Write $L = L_1 \mathbin{+\!+} L_2$ where the insertion happens at the
boundary, so $L' = L_1 \mathbin{+\!+} [\unit] \mathbin{+\!+} L_2$.  Apply
Lemma~\ref{lem:chain-append} twice:
\[
\chain{L'} = \chain{L_1} \compose \chain{[\unit]} \compose \chain{L_2}
= \chain{L_1} \compose \unit \compose \chain{L_2}
= \chain{L_1} \compose \chain{L_2} = \chain{L}.
\]
The last but one equality uses (A3) (and (A2) for the symmetric case
where the inserted unit is at the head).  An inductive argument on the
number of inserted units extends the conclusion to inserting any finite
void chain, recovering Corollary~\ref{cor:1-2} below.  (Verified as
\texttt{IdeaTheory.chain\_concat\_ident} together with
\texttt{chain\_cons\_ident}.)
\end{proof}

\begin{lemma}[Length is additive]\label{lem:chain-length}
$|L_1 \mathbin{+\!+} L_2| = |L_1| + |L_2|$, $|[\,]| = 0$, and
$|a :: L| = |L| + 1$.
\end{lemma}

\begin{proof}
Standard properties of list length, restated for the record.  (Verified as
\texttt{IdeaTheory.chainLength\_append}, \texttt{chainLength\_nil},
\texttt{chainLength\_cons}.)
\end{proof}

\begin{remark}[Chains as discourses]\label{rem:chains-discourse}
The chain construction is the formal correlate of any phenomenon in
which ideas are arranged in a temporal or syntactic sequence.  Hume's
``trains of thought'' (\cite{Hume1739}, Bk.~I, Pt.~I, \S IV) are chains;
so are the cultural attractor sequences of Sperber, the inferential
chains of Brandom, the syntactic compositions of Frege, and the
narrative episodes of Ricoeur.  The bracket-independence guaranteed by
Lemma~\ref{lem:chain-split} is what entitles all of these authors to
treat their primary unit of analysis as a sequence rather than a tree.
\end{remark}

\begin{remark}[Three notions of ``the same'' chain]
The chain construction induces three natural equivalence relations on
finite lists of ideas, of decreasing strictness.  \emph{Literal equality}
demands that the lists agree term by term.  \emph{Composite equality}
identifies $L_1$ and $L_2$ whenever $\chain{L_1} = \chain{L_2}$; by
Lemma~\ref{lem:chain-pad} this collapses any difference produced by
inserting or deleting copies of $\unit$, and by
Lemma~\ref{lem:chain-append} it also collapses bracketing distinctions
already absorbed by (A1).  \emph{Resonance equality}, introduced in
Chapter~2, will further identify chains whose composites differ only by
elements of the radical of $\rho$.  The bare monoid sees only the first
two; the third requires the analytic structure of the next chapter.  This
hierarchy of equivalences is, in our view, the clearest single benefit of
the formal approach: it makes precise the otherwise vague claim that two
discourses ``mean the same thing''.
\end{remark}

\section{Sub-ideals and idempotents}\label{sec:idempotents}

Two derived notions occupy an intermediate position between the bare
algebra and the analytic structure of later chapters.  The first is the
\emph{sub-ideal} relation, which formalizes ``$a$ is contained in $b$''.
The second is the property of \emph{idempotency}, which formalizes ``$a$
absorbs itself''.

\begin{definition}[Left and right sub-ideals]\label{def:subideal}
We say that $a$ is a \emph{left sub-ideal} of $b$, and write
$a \le_L b$, if there exists $c \in \Ideas$ with $b = a \compose c$.
Symmetrically, $a$ is a \emph{right sub-ideal} of $b$, written
$a \le_R b$, if there exists $c$ with $b = c \compose a$.
\end{definition}

\begin{lemma}[Order properties of $\le_L$ and $\le_R$]\label{lem:subideal-order}
The relations $\le_L$ and $\le_R$ are reflexive and transitive.  The void
idea $\unit$ is below every idea: $\unit \le_L a$ and $\unit \le_R a$ for
all $a$.  In particular, $a \le_L a \compose b$ and $b \le_R a \compose b$
for all $a, b$.
\end{lemma}

\begin{proof}
Reflexivity: $a = a \compose \unit$ by (A3), so $a \le_L a$ with witness
$c = \unit$.  Symmetric for $\le_R$ via (A2).  Transitivity: if
$b = a \compose x$ and $c = b \compose y$, then
$c = (a \compose x) \compose y = a \compose (x \compose y)$ by (A1), so
$a \le_L c$ with witness $x \compose y$.  Symmetric for $\le_R$.
$\unit \le_L a$ with witness $c = a$ since $a = \unit \compose a$ by (A2);
symmetric for $\le_R$.  The final two claims are immediate from the
definitions: $a \compose b = a \compose b$ exhibits $a$ as a left factor,
and dually for the right.  (Verified as \texttt{IdeaTheory.subideal\_refl},
\texttt{subideal\_trans}, \texttt{subideal\_ident},
\texttt{subideal\_self\_op\_right}, and the symmetric \texttt{rsubideal\_*}
versions.)
\end{proof}

\begin{remark}
The relations $\le_L$ and $\le_R$ are not in general antisymmetric: in a
group every element is both a left and a right sub-ideal of every other.
Antisymmetry holds in \emph{conical} idea algebras, a special case to
which we shall return in Volume~II.
\end{remark}

\begin{definition}[Idempotent and commuting ideas]\label{def:idempotent}
An idea $a \in \Ideas$ is \emph{idempotent} if $a \compose a = a$.  Two
ideas $a$ and $b$ \emph{commute} if $a \compose b = b \compose a$; we
write $a \mathrel{\#} b$ for this relation.  The \emph{centralizer} of $a$
is the set $Z(a) \defeq \{ b \in \Ideas \mid a \mathrel{\#} b \}$.
\end{definition}

\begin{lemma}[Basic facts about idempotents]\label{lem:idem-basic}
The following hold for every $a \in \Ideas$.
\begin{enumerate}
\item $\unit$ is idempotent.
\item If $a$ is idempotent and $n \ge 1$, then $a^n = a$ (and so $a^n$ is
itself idempotent).
\end{enumerate}
\end{lemma}

\begin{proof}
(1) Lemma~\ref{lem:unit-self}.  (2) Induct on $n$.  Base case $n = 1$:
$a^1 = a$ by Lemma~\ref{lem:npow-elem}.  Inductive step: suppose $a^k = a$
with $k \ge 1$.  Then $a^{k+1} = a \compose a^k = a \compose a = a$ by the
idempotency of $a$.  Idempotency of $a^n$ follows since $a^n \compose a^n
= a \compose a = a = a^n$.  (Verified as
\texttt{IdeaTheory.ident\_idempotent} and
\texttt{IdeaTheory.npow\_idempotent}, \texttt{idempotent\_npow}.)
\end{proof}

\begin{lemma}[Basic facts about commutation]\label{lem:commute-basic}
The relation $\mathrel{\#}$ is reflexive and symmetric; the void idea
commutes with everything; in particular $\unit \in Z(a)$ and $a \in Z(a)$
for every $a$.
\end{lemma}

\begin{proof}
Reflexivity is immediate.  Symmetry: if $a \compose b = b \compose a$,
then equality is symmetric, so $b \compose a = a \compose b$.
$\unit$ commutes with $a$: $\unit \compose a = a = a \compose \unit$ by
(A2)--(A3).  The membership claims follow.  (Verified as
\texttt{IdeaTheory.commute\_refl}, \texttt{commute\_symm},
\texttt{commute\_ident\_left}, \texttt{commute\_ident\_right},
\texttt{centralizer\_ident}, \texttt{centralizer\_self}.)
\end{proof}

The relation $\mathrel{\#}$ is \emph{not} in general transitive: $a$ and
$b$ may both commute with $c$ without commuting with each other.  This
failure is one of the algebraic shadows of the non-commutativity of
$\compose$ and reappears in Chapter~5 as the curvature of the resonance
pairing.

\section{Cancellation and invertibility}\label{sec:cancellation}

We now turn to two properties that distinguish ``rich'' ideas (whose
content is sufficient to be cancelled in equations or even reversed) from
generic ones.

\begin{definition}[Cancellative ideas]\label{def:cancellative}
An idea $a \in \Ideas$ is
\begin{itemize}
\item \emph{left-cancellative} if for all $b, c$, the equation
$a \compose b = a \compose c$ implies $b = c$;
\item \emph{right-cancellative} if for all $b, c$, the equation
$b \compose a = c \compose a$ implies $b = c$;
\item \emph{cancellative} if both.
\end{itemize}
\end{definition}

\begin{lemma}[$\unit$ is cancellative]\label{lem:unit-cancel}
$\unit$ is both left- and right-cancellative.
\end{lemma}

\begin{proof}
If $\unit \compose b = \unit \compose c$, then by (A2) $b = c$.
Symmetric on the right.  (Verified as \texttt{IdeaTheory.ident\_left\_cancel},
\texttt{ident\_right\_cancel}, \texttt{ident\_cancellative}.)
\end{proof}

\begin{lemma}[Composition of cancellative ideas]\label{lem:cancel-compose}
If $a$ and $b$ are left-cancellative, so is $a \compose b$.  Symmetrically
for right cancellation, and hence for two-sided cancellation.
\end{lemma}

\begin{proof}
Suppose $(a \compose b) \compose x = (a \compose b) \compose y$.  By (A1)
this is $a \compose (b \compose x) = a \compose (b \compose y)$.  Left
cancellation by $a$ gives $b \compose x = b \compose y$, then left
cancellation by $b$ gives $x = y$.  The dual argument handles the right
case.  (Verified as \texttt{IdeaTheory.left\_cancel\_op},
\texttt{right\_cancel\_op}.)
\end{proof}

The property of \emph{invertibility} is sharper still: an idea is
invertible if it has a two-sided ``undoing'' --- an inverse that, composed
on either side, returns the void.

\begin{definition}[Invertible ideas]\label{def:invertible}
An idea $a \in \Ideas$ is \emph{invertible} if there exists $b \in \Ideas$
with $a \compose b = \unit$ and $b \compose a = \unit$.  In that case $b$
is called an \emph{inverse} of $a$.
\end{definition}

\begin{lemma}[$\unit$ is invertible]\label{lem:unit-inv}
$\unit$ is invertible, with inverse $\unit$.
\end{lemma}

\begin{proof}
$\unit \compose \unit = \unit$ by Lemma~\ref{lem:unit-self}.
(Verified as \texttt{IdeaTheory.ident\_invertible}.)
\end{proof}

\begin{lemma}[Inverses are unique]\label{lem:inv-unique}
If $b_1$ and $b_2$ are both inverses of $a$, then $b_1 = b_2$.
\end{lemma}

\begin{proof}
We have $a \compose b_2 = \unit$ and $b_1 \compose a = \unit$.  Compute
\[
b_1 \;=\; b_1 \compose \unit
     \;=\; b_1 \compose (a \compose b_2)
     \;=\; (b_1 \compose a) \compose b_2
     \;=\; \unit \compose b_2
     \;=\; b_2,
\]
using (A3), the right-inverse property of $b_2$, (A1), the left-inverse
property of $b_1$, and (A2) in turn.  (Verified as
\texttt{IdeaTheory.invertible\_inv\_unique}.)
\end{proof}

\begin{lemma}[Invertible ideas are closed under composition]\label{lem:inv-compose}
If $a$ and $b$ are invertible, with inverses $a'$ and $b'$ respectively,
then $a \compose b$ is invertible with inverse $b' \compose a'$.
\end{lemma}

\begin{proof}
Compute the two required identities.  For the right-inverse property,
\[
(a \compose b) \compose (b' \compose a')
= a \compose (b \compose (b' \compose a'))
= a \compose ((b \compose b') \compose a')
= a \compose (\unit \compose a')
= a \compose a'
= \unit,
\]
using (A1) twice, the right-inverse property of $b'$, (A2), and the
right-inverse property of $a'$.  The left-inverse property is symmetric:
\[
(b' \compose a') \compose (a \compose b)
= b' \compose ((a' \compose a) \compose b)
= b' \compose (\unit \compose b)
= b' \compose b
= \unit.
\]
(Verified as \texttt{IdeaTheory.invertible\_op}.)
\end{proof}

The set $\Ideas^{\times} \defeq \{ a \in \Ideas : a \text{ is invertible}\}$
is therefore closed under composition and contains $\unit$; it is, by
construction, a \emph{group}.  We single this out as a definition for
later reference.

\begin{definition}[The group of units]\label{def:units}
The \emph{group of units} of an idea algebra $\IdMon$ is the subset
$\Ideas^{\times}$ of invertible ideas, equipped with the restriction of
$\compose$, the unit $\unit$, and the inversion map $a \mapsto a^{-1}$
(well defined by Lemma~\ref{lem:inv-unique}).
\end{definition}

\begin{remark}
The group $\Ideas^{\times}$ measures the \emph{reversibility budget} of
the algebra.  In a free monoid (Locke's atomistic picture, see
\S\ref{sec:zoo}) we shall see that $\Ideas^{\times} = \{\unit\}$: nothing
is reversible.  At the opposite extreme, in a group every idea is
invertible.  Realistic models of cognition lie in between, and the size of
$\Ideas^{\times}$ is then a measurable structural invariant.
\end{remark}

\begin{table}[h]
\centering
\begin{tabular}{ll}
\hline
Symbol & Meaning \\
\hline
$\compose$ & Composition of ideas \\
$\unit$ & Void idea (unit) \\
$a^n$ & Iterated composition $\npow(n, a)$ \\
$\chain{L}$ & Composition of a list of ideas \\
$|L|$ & Length of a chain \\
$a \le_L b$ / $a \le_R b$ & Left / right sub-ideal \\
$a \mathrel{\#} b$ & Commutation \\
$Z(a)$ & Centralizer of $a$ \\
$\Ideas^{\times}$ & Group of invertible ideas \\
\hline
\end{tabular}
\caption{Notation introduced in Chapter~\ref{ch:monoid}.}
\label{tab:notation-ch1}
\end{table}

\section{Theorems 1.1--1.3 and corollaries 1.1--1.4}\label{sec:headline}

We now consolidate the structural payload of the chapter into three
theorems and four corollaries.  All seven results are formally verified
in the Lean module \texttt{IdeaTheory.Theorems1}; we cite the exact
identifier in each proof.

\begin{theorem}[Bracketing-independence of chain composition]\label{thm:1-1}
For every finite list $L$ of ideas and every $n \in \NN$,
\begin{equation}\label{eq:thm-1-1}
\chain{L}
\;=\; \chain{\mathrm{take}_n L} \compose \chain{\mathrm{drop}_n L}.
\end{equation}
\end{theorem}

\begin{proof}
Write $L = L_1 \mathbin{+\!+} L_2$ with $L_1 = \mathrm{take}_n L$ and
$L_2 = \mathrm{drop}_n L$.  Apply Lemma~\ref{lem:chain-append} to obtain
$\chain{L_1 \mathbin{+\!+} L_2} = \chain{L_1} \compose \chain{L_2}$, which
is the right-hand side of \eqref{eq:thm-1-1}.  Substituting back yields
the left-hand side.  (Verified as \texttt{IdeaTheory.theorem\_1\_1}.)
\end{proof}

Theorem~\ref{thm:1-1} is the algebraic guarantee that the composite
content of a discourse, narrative, or transmission chain depends only on
the linear order of its parts and not on any prior decision about how to
bracket them.  This is the formal version of the informal principle,
common to authors as different as \cite{Hume1739} and Sperber, that ``the
order in which we group the units does not affect the result''.

\begin{theorem}[Uniqueness of the void; closure of the units; uniqueness of inverses]\label{thm:1-2}
Let $\IdMon$ be a bare idea algebra.  Then:
\begin{enumerate}
\item If $e \in \Ideas$ is a two-sided unit, then $e = \unit$.
\item If $a, b \in \Ideas$ are invertible, then so is $a \compose b$.
\item If $b_1, b_2$ are inverses of the same idea $a$, then $b_1 = b_2$.
\end{enumerate}
\end{theorem}

\begin{proof}
(1) Lemma~\ref{lem:unit-unique-twosided}.  (2)
Lemma~\ref{lem:inv-compose}.  (3) Lemma~\ref{lem:inv-unique}.  Each
clause has been proved separately above; bundling them is a matter of
notation.  (Verified as \texttt{IdeaTheory.theorem\_1\_2}.)
\end{proof}

\begin{theorem}[Iteration is a monoid homomorphism]\label{thm:1-3}
For every fixed $a \in \Ideas$ the function $n \mapsto a^n$ from $\NN$ to
$\Ideas$ satisfies
\begin{equation}\label{eq:thm-1-3-zero}
a^0 = \unit
\end{equation}
and
\begin{equation}\label{eq:thm-1-3-add}
a^{m+n} = a^m \compose a^n,
\quad
a^{m+n+k} = (a^m \compose a^n) \compose a^k,
\end{equation}
for all $m, n, k \in \NN$.  Equivalently, $n \mapsto a^n$ is a monoid
homomorphism $(\NN, +, 0) \to \IdMon$.
\end{theorem}

\begin{proof}
Equation~\eqref{eq:thm-1-3-zero} is Definition~\ref{def:npow}.  The first
identity in \eqref{eq:thm-1-3-add} is Lemma~\ref{lem:npow-add}.  For the
second, apply Lemma~\ref{lem:npow-add} twice:
\[
a^{m+n+k} = a^{(m+n)} \compose a^k
        = (a^m \compose a^n) \compose a^k.
\]
Together with Lemma~\ref{lem:unit-self}, these identities are exactly the
defining properties of a monoid homomorphism.  (Verified as
\texttt{IdeaTheory.theorem\_1\_3}.)
\end{proof}

\begin{remark}[Why we package these specific results]
Theorems \ref{thm:1-1}, \ref{thm:1-2}, \ref{thm:1-3} are not the only
non-trivial facts derivable from the axioms; they are the three that are
\emph{used} in every later chapter.  Theorem~\ref{thm:1-1} licenses
manipulation of chain composites; Theorem~\ref{thm:1-2} guarantees that
the unit and inverses behave canonically; Theorem~\ref{thm:1-3} is the
input to every later argument involving iteration of an idea, and in
particular feeds directly into the cocycle calculations of Chapter~5.
\end{remark}

The four corollaries below connect the three theorems to specific
applications taken up in later volumes.

\begin{corollary}[Reorganization of cultural transmission chains]\label{cor:1-1}
For any finite lists $L_1, L_2, L_3$ of ideas,
\begin{equation}\label{eq:cor-1-1}
\chain{L_1 \mathbin{+\!+} L_2 \mathbin{+\!+} L_3}
\;=\; \chain{L_1 \mathbin{+\!+} (L_2 \mathbin{+\!+} L_3)}.
\end{equation}
\end{corollary}

\begin{proof}
List concatenation is associative on the nose: $L_1 \mathbin{+\!+} L_2
\mathbin{+\!+} L_3 = L_1 \mathbin{+\!+} (L_2 \mathbin{+\!+} L_3)$.  Apply
$\chain{\,\cdot\,}$.  (Verified as \texttt{IdeaTheory.corollary\_1\_1}.)
This corollary will be invoked in Volume~II to prove that diffusion of an
idea through a population is independent of the choice of intermediate
broadcast points.
\end{proof}

\begin{corollary}[Insertion of empty episodes]\label{cor:1-2}
For any finite lists $L_1, L_2$ and any $n \in \NN$,
\begin{equation}\label{eq:cor-1-2}
\chain{L_1 \mathbin{+\!+} \unit^{::n} \mathbin{+\!+} L_2}
\;=\; \chain{L_1 \mathbin{+\!+} L_2}.
\end{equation}
\end{corollary}

\begin{proof}
By Lemma~\ref{lem:chain-append} applied twice and
Lemma~\ref{lem:chain-basic}\eqref{eq:chain-void},
\begin{align*}
\chain{L_1 \mathbin{+\!+} \unit^{::n} \mathbin{+\!+} L_2}
&= \chain{L_1 \mathbin{+\!+} \unit^{::n}} \compose \chain{L_2} \\
&= (\chain{L_1} \compose \unit) \compose \chain{L_2}
= \chain{L_1} \compose \chain{L_2}
= \chain{L_1 \mathbin{+\!+} L_2},
\end{align*}
using (A3) at the third step.  (Verified as
\texttt{IdeaTheory.corollary\_1\_2}.)  The corollary will be used in
Volume~III to prove that two narratives differing only by the insertion
of empty interludes have the same total content.
\end{proof}

\begin{corollary}[Closure of reversible operations]\label{cor:1-3}
The set $\Ideas^{\times}$ of invertible ideas is closed under
composition.  In particular, if $a$ and $b$ are invertible, so is
$a \compose b$.
\end{corollary}

\begin{proof}
Lemma~\ref{lem:inv-compose}, repackaged.  (Verified as
\texttt{IdeaTheory.corollary\_1\_3}.)  The intended application in
Volume~IV is to mental operations equipped with an undo: that the
composition of two undoable operations is itself undoable is what
licenses backtracking and counterfactual reasoning.
\end{proof}

\begin{corollary}[Idempotents are exhausted by themselves]\label{cor:1-4}
If $a \in \Ideas$ is idempotent, then $a^n = a$ for every $n \ge 1$.
\end{corollary}

\begin{proof}
Lemma~\ref{lem:idem-basic}(2).  (Verified as
\texttt{IdeaTheory.corollary\_1\_4}.)  The corollary expresses an
emergence-theoretic fact taken up in Volume~V: an absorbed cultural item
does not grow under repetition; whatever surplus emergence might have
contributed has already been delivered after a single application.
\end{proof}

\section{Three concrete examples}\label{sec:examples}

Before turning to the historical zoo, we exhibit three concrete idea
algebras to ground the formalism.  Each is verified in the Lean module as
an instance of \texttt{IdeaTheoryStructure}; here we describe them
mathematically.

\begin{example}[The natural numbers under addition]\label{ex:nat}
Take $\Ideas = \NN$, $\compose = +$, and $\unit = 0$.  Axioms (A1)--(A3)
are the standard properties of addition.  In this model a chain of
``ideas'' $[n_1, \dots, n_k]$ has composite $n_1 + \cdots + n_k$, and the
$n$-th power of $a$ is $na$.  Idempotents are exhausted by $0$;
invertible elements are again exhausted by $0$.  This trivial-looking
example is in fact a sharp test case: it shows that the bare monoid axioms
are consistent and that none of (A1)--(A3) is vacuous (each axiom can be
exhibited independently in this model).  More substantively, the model
demonstrates that the formalism does not by itself force any non-trivial
group of units --- $\Ideas^{\times}$ may consist of $\unit$ alone.
\end{example}

\begin{example}[Free idea algebra on a generating set $S$]\label{ex:free}
Let $S$ be any non-empty set and let $\Ideas = S^{*}$ denote the set of
finite words on $S$.  Define $\compose$ as concatenation and $\unit$ as
the empty word.  This is the free monoid on $S$, and it is the
fundamental object underlying the Lockean reading of
\S\ref{zoo:locke}.  Two facts about this model are worth stating
explicitly.  First, the chain composite of a list $[w_1, \dots, w_k]$ of
words is the concatenation $w_1 w_2 \cdots w_k$, recovering the intuition
that ``adding more ideas'' literally lengthens the resulting expression.
Second, the only idempotent is $\unit$, and the only invertible element
is again $\unit$: in the free idea algebra nothing absorbs and nothing is
reversible.  These properties are the algebraic shadow of Locke's
``inability to invent a new simple idea'': there is no internal mechanism
for the algebra to collapse two distinct combinations of generators into
the same element.
\end{example}

\begin{example}[A two-atom finite world]\label{ex:two-atom}
Let $\Ideas = \{ \unit, a, b, c \}$ with multiplication table given by
$a \compose a = a$, $b \compose b = b$, $a \compose b = b \compose a = c$,
$c \compose x = c$ for $x \in \{a, b, c\}$, and $\unit$ acting as the
two-sided unit.  One verifies directly that this is associative and so
constitutes a bare idea algebra.  Both $a$ and $b$ are idempotent;
$c$ is the ``saturated'' idea that absorbs everything.  This toy model
has no non-trivial invertibles, but it has multiple idempotents, making it
a useful playground for the lemmas of \S\ref{sec:idempotents}.  The model
also illustrates how the bare monoid can encode the informal idea of
``saturation by absorption'': once $c$ is reached, no further composition
moves the system.  In the language of Volume~V, $c$ is a fixed point of
the cultural-attractor dynamics.
\end{example}

These three models bracket the spectrum of bare idea algebras: the
arithmetic model (commutative, no non-trivial idempotents, no
non-trivial invertibles), the free model (non-commutative, no
idempotents, no invertibles), and the saturation model (non-commutative,
many idempotents, an absorbing element).  Every later chapter will return
to these examples to illustrate the additional structure being introduced.

\begin{remark}[On the role of examples]
The role of examples in a structural theory is to certify that the
abstract definitions are not accidentally inconsistent and that the
theorems we prove from them have non-trivial content.  The three examples
above suffice for both purposes.  Example~\ref{ex:nat} shows
Theorem~\ref{thm:1-3} is non-trivial (the multiplicative law $a^{m+n} =
a^m + a^n$ becomes the familiar identity $(m+n)a = ma + na$).
Example~\ref{ex:free} shows that Theorem~\ref{thm:1-1} is
non-trivial (different positions of split give different intermediate
words, but their concatenation always reconstructs the original).
Example~\ref{ex:two-atom} shows that Corollary~\ref{cor:1-4} has bite
($c$ is idempotent and indeed $c^n = c$ for all $n \ge 1$).
\end{remark}

\begin{remark}[A note on cardinality]
The bare monoid axioms place no constraints on the cardinality of
$\Ideas$.  All three of our examples have countable carriers:
Example~\ref{ex:nat} has $|\Ideas| = \aleph_0$, Example~\ref{ex:free}
has $|\Ideas| = \aleph_0$ when $S$ is countable, and Example~\ref{ex:two-atom}
has $|\Ideas| = 4$.  Realistic models of natural-language semantics or
of cultural transmission typically demand uncountable carriers (often
the underlying set is, or contains, a topological space such as a
function space or a manifold).  None of the results of this chapter
depend on cardinality assumptions, and we shall freely apply them to
uncountable models when the need arises in later volumes.  The cost of
this generality is paid in the analytic structure of Chapter~2, where
the resonance pairing must be defined coherently on infinite-dimensional
free vector spaces.
\end{remark}

\section{The zoo: four classical positions accommodated by the bare monoid}\label{sec:zoo}

We now show that the bare monoid structure already suffices to host four
canonical positions in the philosophy of ideas: the atomistic combinatorics
of \cite{Locke1690}, the associationist chains of \cite{Hume1739}, the
sense-composition of \cite{Frege1892}, and the family-resemblance picture
of the later \cite{Wittgenstein1953}.  Each entry follows the
brief-mandated template (\emph{Original Claim} / \emph{Formal Restatement}
/ \emph{Status} / \emph{Commentary}) and contains at least two citations.

\subsection*{Locke (1690): simple and complex ideas as monoid generators}\label{zoo:locke}

\paragraph{Original claim.}
Locke divides ideas into the \emph{simple} and the \emph{complex}
(\cite{Locke1690}, Bk.~II, Ch.~II--XII).  Simple ideas are received
passively from sensation or reflection; the mind cannot manufacture them.
Complex ideas, by contrast, are produced by acts of \emph{combination},
\emph{comparison}, and \emph{abstraction} performed on simple ideas (and
on previously formed complex ideas).  The mind's productive power is
exhausted by these operations: ``it being not in the power of the most
exalted wit, or enlarged understanding, by any quickness or variety of
thought, to invent or frame one new simple idea in the mind'' (Bk.~II,
Ch.~II, \S 2).  Locke is here explicit about a closure principle: every
idea in the mind is either simple or generated from simples by finitely
many combinations.  See also \cite{Hume1739}, Bk.~I, Pt.~I, \S I, where
Hume's distinction between impressions and ideas inherits Locke's
combinatorial picture and sharpens the closure principle into the
``copy principle''.

\paragraph{Formal restatement.}
Let $S \subseteq \Ideas$ be the set of simple ideas.  Locke's claim is
that $\Ideas$ is the smallest sub-monoid of itself containing $S$ and
closed under $\compose$.  Equivalently, the inclusion $S \hookrightarrow
\Ideas$ extends to a unique surjective monoid homomorphism from the free
monoid $S^{*}$ onto $\Ideas$.  Concretely, every $a \in \Ideas$ is of the
form $\chain{[s_1, \dots, s_k]}$ for some $s_1, \dots, s_k \in S \cup
\{\unit\}$ and some $k \ge 0$.

\paragraph{Status.}
\textbf{True for any choice of generators $S$.}  By
Lemma~\ref{lem:chain-append}, the set of chain composites of finite
sequences from $S \cup \{\unit\}$ is closed under $\compose$ and contains
$\unit = \chain{[\,]}$, hence is a sub-monoid; and Lemma~\ref{lem:chain-basic}
shows that it contains every singleton.  The non-trivial content of
Locke's claim --- that $\Ideas$ \emph{equals} the smallest such sub-monoid
for some philosophically meaningful $S$ --- is then a question about the
existence of a generating set, not about the algebra.

\paragraph{Commentary.}
The bare monoid is faithful to Locke at the structural level but silent on
two points where Locke himself is silent: it does not say \emph{which}
ideas count as simple, and it does not predict whether the generation map
$S^{*} \to \Ideas$ is injective (i.e.\ whether different bracketed
combinations give different complex ideas).  The latter question is
addressed in Volume~I, Chapter~4, where the kernel of the generation map
is computed in terms of the resonance pairing.  An empirically charged
prediction follows: in a model where $S^{*} \to \Ideas$ is injective,
Lockean atomism is exactly correct; in a model where it is not, two
distinct combinations of simples can produce the same complex idea --- the
algebraic shadow of \emph{Aufhebung}, taken up systematically in
Chapter~3.  A second point of departure from Locke is the status of
\emph{abstraction}: Locke treats abstraction as a third primitive
operation alongside combination and comparison, but the formalism
recovers abstraction as the passage from $\Ideas$ to a quotient
sub-monoid (cf.\ \S\ref{zoo:witt}), making it derivative rather than
primitive.  This is a substantive divergence whose empirical content will
be tested in Volume~IV against developmental data on concept formation.

\subsection*{Hume (1739): chains as principles of association}\label{zoo:hume}

\paragraph{Original claim.}
\cite{Hume1739} (Bk.~I, Pt.~I, \S IV) holds that ideas combine in the
mind by three universal principles of association: resemblance,
contiguity in time and place, and cause and effect.  These principles
operate on \emph{trains} of ideas --- sequences in which each new term is
recruited by association with its predecessor.  The content of any
complex thought is then the result of running such a train of associations
to its end.  Hume's claim is doubly strong: first, that all complex
thinking is sequential association; second, that the principles
producing the sequence are exhausted by the three he names.  Hume's
treatment is sharpened by \cite{Locke1690} (Bk.~II, Ch.~XXXIII), whose
own discussion of the ``association of ideas'' is Hume's immediate source.

\paragraph{Formal restatement.}
Identify a Humean train of ideas with a finite list $L = [a_1, \dots,
a_k]$ of elements of $\Ideas$.  The content of the train is its chain
composite $\chain{L}$.  Hume's first claim, in this notation, is that
every complex idea is of this form.  Hume's second claim --- that the
sequence of $a_i$ is generated by a fixed associative process --- is the
demand that there be a (partial) function $\alpha : \Ideas \to \Ideas$
with $a_{i+1} = \alpha(a_i)$ for every $i$, where $\alpha$ encodes the
three association principles.

\paragraph{Status.}
\textbf{Hume's first claim is true on the bare monoid.}  By
Lemma~\ref{lem:chain-basic} every idea $a$ is itself a chain
($\chain{[a]} = a$), and by Lemma~\ref{lem:chain-append} every composite
of a finite collection of ideas is a chain.  The claim that complex
ideas \emph{are} chains is therefore vacuous given (A1)--(A3).
\textbf{Hume's second claim is contingent.}  It can be made to hold or
fail by appropriate choice of $\alpha$ on the underlying carrier $\Ideas$.
The bare monoid is silent on the dynamics; this silence is a feature, not
a bug, since later volumes equip $\Ideas$ with the resonance pairing
$\rho$ that makes ``association by resemblance'' formally tractable as a
proximity in the resonance metric.

\paragraph{Commentary.}
Hume's chain principle is one of the most under-defended assumptions in
the early-modern theory of mind, and the formalization makes clear why:
on the bare monoid it is mathematically empty.  The substantive content
of associationism only becomes visible once one fixes (i) a generating
set and (ii) a dynamics on it; both are deferred to later chapters.
The empirical prediction unique to the formal version is
Lemma~\ref{lem:chain-pad}: empty episodes inserted into a Humean train
do not change its content.  This prediction is testable: any
psychological theory in which inserting a contentless interval changes
the train's outcome is incompatible with axioms (A1)--(A3) and so
demands a richer underlying algebra.

\subsection*{Frege (1892): composition of senses}\label{zoo:frege}

\paragraph{Original claim.}
\cite{Frege1892} introduces the distinction between \emph{Sinn} (sense)
and \emph{Bedeutung} (reference) for linguistic expressions.  The sense of
a complex expression is determined by the senses of its parts and their
mode of combination --- the principle of \emph{compositionality}.  Frege's
own examples (the morning star and the evening star, the king of France)
emphasize that two expressions can share their reference without sharing
their sense, and that sense composes systematically as one builds up
complex expressions from simple ones.  The principle is reaffirmed in
\cite{Wittgenstein1922}, \S 3.318, which treats the proposition as a
function of its constituent expressions.

\paragraph{Formal restatement.}
Let $\Ideas$ be the set of senses, and let $\compose$ be the operation
``the sense of $a$-applied-to-$b$'' (in Frege's setting, the saturation
of a function-sense by an argument-sense).  Compositionality is the
demand that $\compose$ be a \emph{function} of its inputs: the sense of
the whole is fully determined by the senses of the parts.  Together with
the natural unit (the sense of the empty context), this yields exactly
the bare idea algebra $\IdMon$.

\paragraph{Status.}
\textbf{True, but the substantive content is associativity.}  Frege's
compositionality principle is, on its face, only the assertion that
$\compose$ is well defined as a function.  Axiom (A1) is what makes the
principle useful: it guarantees that the sense of a complex expression
does not depend on the order in which its sub-expressions are saturated.
Without (A1), Fregean composition would collapse into a tree-structured
operation in which every re-bracketing produces a different sense.

\paragraph{Commentary.}
The bare monoid is faithful to Frege provided one is willing to read
Frege's notion of ``mode of combination'' as itself encoded inside the
operation $\compose$ (rather than as additional structure).  The
formalization sharpens a tension in Frege's text: Frege's own
\emph{Function and Concept} treats functional application as
non-commutative, and the formal apparatus inherits this and is
philosophically faithful to it.  A genuine prediction of the formalization
that goes beyond Frege is Theorem~\ref{thm:1-1}: the sense of a complex
expression is invariant under any contiguous re-grouping of its
constituents, not merely under binary re-association.  The role of the
emergence cocycle $\omega$ (Chapter~5) is to measure exactly the
non-Fregean residue of composition --- the portion of meaning that is
\emph{not} a function of the senses of the parts.

\subsection*{Wittgenstein (1953): family resemblance as a quotient monoid}\label{zoo:witt}

\paragraph{Original claim.}
The later \cite{Wittgenstein1953} (\S 65--67) rejects the Lockean--Fregean
picture of concepts as having essential content.  Instead, the extension
of a concept --- Wittgenstein's example is ``game'' --- is held together
by a network of overlapping similarities, ``a complicated network of
similarities overlapping and criss-crossing'', which Wittgenstein labels
\emph{family resemblance} (\emph{Familien\"ahnlichkeit}).  Members of a
family-resemblance class need share no single feature; what binds them is
that any two are connected by a chain of pairwise similarities.  The
position is anticipated by Hume's discussion of resemblance as one of the
three principles of association (\cite{Hume1739}, Bk.~I, Pt.~I, \S IV).

\paragraph{Formal restatement.}
Fix an equivalence relation $\equivIdea$ on $\Ideas$ generated by a
symmetric ``resemblance'' relation $\sim$.  The family-resemblance class
of $a$ is then $[a] \defeq \{ b : a \equivIdea b \}$.  Wittgenstein's
claim is that the natural object of conceptual analysis is not $\Ideas$
but the quotient $\Ideas / {\equivIdea}$.  The bare monoid hosts this
construction provided $\equivIdea$ is a \emph{congruence} for $\compose$:
that is, $a_1 \equivIdea a_2$ and $b_1 \equivIdea b_2$ imply
$a_1 \compose b_1 \equivIdea a_2 \compose b_2$.

\paragraph{Status.}
\textbf{Contingent on $\equivIdea$ being a congruence.}  When this holds,
the quotient $\Ideas / {\equivIdea}$ inherits a well-defined operation
$[a] \compose [b] \defeq [a \compose b]$ and the quotient $(\Ideas /
{\equivIdea}, \compose, [\unit])$ is itself a bare idea algebra; the
quotient map is a surjective homomorphism.  This is the universal
property of quotient monoids (\cite{MacLane1998}, Ch.~I, \S 4;
\cite{Bourbaki1989}, Ch.~I, \S 4, no.~9).  When $\equivIdea$ is not a
congruence --- which is the typical case for genuine family resemblance,
since composing two ``games'' need not yield something resembling a
``game'' --- the bare monoid structure does \emph{not} descend to the
quotient.  In that case Wittgenstein's class is not closed under any
algebraic operation derived from $\compose$, and what survives at the
quotient is only a relation of resemblance, not a monoid.

\paragraph{Commentary.}
The formalization sharply discriminates two readings of family
resemblance.  On the \emph{algebraic} reading, family resemblance is a
congruence and Wittgenstein's classes are quotient ideas: this is
faithful to those passages where Wittgenstein speaks of concepts as
determining patterns of use.  On the \emph{relational} reading, family
resemblance is a mere relation that need not respect composition: this
is faithful to those passages where Wittgenstein insists that the
boundaries of a concept are open and inferentially uncontrolled.  The
choice between the two readings is forced once a candidate $\equivIdea$
is specified, by the testable question of whether it is a congruence.
The bare monoid is what makes the question testable.  This is exactly the
diagnostic role we expect the formalism to play throughout the zoo: not
to settle exegetical disputes by fiat but to make the empirical
consequences of each reading explicit.

\medskip

\noindent
With these four entries the bare monoid has accommodated, without further
structure, a recognizably wide swath of the classical literature on
ideas.  Subsequent chapters add the resonance pairing
(Chapter~\ref{ch:monoid}+1), the Aufhebung decomposition, the emergence
cocycle, and finally the grading; each of these layers will host further
zoo entries, and each layer remains anchored to the present chapter
through the underlying monoid axioms (A1)--(A3).

\section{Synopsis and what comes next}\label{sec:synopsis}

We close the chapter by stepping back from the technical apparatus and
recording, in a single page, what has been accomplished and what remains
to be done.

\paragraph{Accomplishments.}
We have isolated three axioms (A1)--(A3) and shown that they suffice to
derive: the uniqueness of the void idea (Lemma~\ref{lem:unit-unique-twosided});
the basic arithmetic of iteration $a^n$ (Lemmas~\ref{lem:npow-elem}--%
\ref{lem:npow-mul}); the homomorphism property of chain composition
(Lemmas~\ref{lem:chain-append}--\ref{lem:chain-pad}); the order
properties of the sub-ideal relations (Lemma~\ref{lem:subideal-order});
the closure of idempotents and commuting pairs under the appropriate
constructions (Lemmas~\ref{lem:idem-basic}--\ref{lem:commute-basic}); the
uniqueness of inverses and closure of $\Ideas^{\times}$ under composition
(Lemmas~\ref{lem:inv-unique}--\ref{lem:inv-compose}); and the three
headline theorems (\ref{thm:1-1}--\ref{thm:1-3}) with their four
corollaries (\ref{cor:1-1}--\ref{cor:1-4}).  Every result is verified in
the Lean module \texttt{IdeaTheory.Theorems1}.

\paragraph{Limitations of the bare monoid.}
The structure $\IdMon$ is silent on three questions of immediate
philosophical interest.  First, it cannot distinguish two ideas that
behave identically under composition; the algebra has no internal
mechanism for ``intentional'' or ``content-level'' difference.  Second,
it provides no measure of \emph{distance} or \emph{similarity} between
ideas; the family-resemblance reading of \cite{Wittgenstein1953} has to
be encoded by an externally specified equivalence
(\S\ref{zoo:witt}).  Third, it has no notion of \emph{weight} or
\emph{magnitude}: a chain consisting of a single dominant idea and many
trivial ones is treated identically with a chain whose ideas all
contribute equally.  Each of these limitations is an opening for the
chapters to come.

\paragraph{Roadmap.}
The next chapter introduces the resonance pairing $\rho$ that supplies
the missing notion of distance, and shows that it extends bilinearly to
the free vector space on $\Ideas$.  Chapter~3 develops the Aufhebung
decomposition $a \compose b = \sigma(a,b) + \nu(a,b) + \eta(a,b)$ and
proves its uniqueness modulo $\ker \rho$.  Chapter~5 introduces the
emergence cocycle $\omega$ that measures the failure of any
content-valuation $v : \Ideas \to A$ to be a monoid homomorphism, and
shows that its cohomology class is a structural invariant of $\IdMon$.
Chapters 6--9 add conjugation, transmission chains, structural
equivalence, and grading respectively.  Each of these chapters will
re-prove, on its richer carrier, an analogue of one of the headline
results of the present chapter; in every case the proof reduces, after
the new structure has been peeled away, to one of the lemmas established
here.

\paragraph{The zoo, continued.}
The four zoo entries of \S\ref{sec:zoo} were chosen to demonstrate that
the bare monoid is already philosophically expressive.  In subsequent
chapters the same template (\emph{Original Claim} / \emph{Formal
Restatement} / \emph{Status} / \emph{Commentary}) will host entries on
Plato (Forms as fixed points of composition; Chapter~3), Hegel
(\emph{Aufhebung} as the elevation component $\eta$; Chapter~4), Kant
(categories as schematic operators acting by conjugation; Chapter~6),
Husserl (noemata as elements of the resonance dual; Chapter~2),
Deleuze (virtual problems as cohomology classes; Chapter~5), Quine
(translation indeterminacy as non-uniqueness of the structural
isomorphism; Chapter~8), Fodor (LOT atoms as monoid generators of a
graded algebra; Chapter~9), Sperber (cultural attractors as fixed
points of $\compose$ followed by truncation in the resonance norm;
Chapter~7), and many others, including the modern cognitive-science
literature on prototypes, frames, conceptual spaces, and conceptual
blending.  In each case, the bare monoid of the present chapter remains
the algebraic skeleton on which the additional structure hangs, and the
diagnostic questions raised by each zoo entry are formulated in the
language of monoid homomorphisms, congruences, and quotients exactly as
they are in \S\ref{zoo:witt}.

\paragraph{A methodological closing remark.}
The conviction underlying the entire monograph is that the abundance of
informal claims about ``what ideas are'' and ``how they combine'' has
created, over centuries, a substantial body of \emph{tacit algebra}.
Untangling this tacit algebra and pinning down its minimal axiomatic
core is the task of Idea Theory.  The bare monoid of this chapter is the
first unit of that effort.  It is small, but it is non-trivial: every
position in the zoo had to be \emph{interpreted}, not merely cited, and
the interpretation in each case made testable predictions that go
beyond the original informal claim.  The next eight chapters will repeat
this exercise on richer structure, and Volume~II will repeat it on
sociological and cultural data.  Throughout, the discipline imposed by
working in a fully verified setting --- the Lean module
\texttt{IdeaTheory} contains no \texttt{sorry} and no \texttt{admit}
--- guarantees that the formal content of every claim is exactly what
the prose says it is.

